%======================================================================
\chapter{The space of elements is a plane}\label{ch:quotient}
%======================================================================

We return to a Moore decomposition $G$ of the plane (Definition~\ref{def:moore-decomp}) and finish
Moore's 1925 paper: Theorems~13--22. The plan of the paper is simple. Theorems~1--12 (Chapters
\ref{ch:usc}--\ref{ch:jordan}) show that arcs and simple closed curves of elements behave as in the plane,
with the Jordan curve theorem for curves of elements (Theorem~\ref{thm:m25-11}) as the key result.
Theorems~13--19 develop the cross-cut theory of curves of elements, which shows that the
interiors of simple closed curves of elements, the \emph{regions in the restricted sense}, form a basis of the
topology. Theorem~21 checks that with these regions $G$ satisfies Moore's Axioms~1--8, and Theorem~22 then
follows from Moore's characterization of the plane (Theorem~\ref{thm:moore-plane}).

Throughout, $J$ denotes a simple closed curve of elements; by Theorem~\ref{thm:m25-11},
$G\setminus J$ has exactly two components, the bounded \emph{interior} $I(J)$ and the unbounded
\emph{exterior} $E(J)$, and $J$ is the frontier of each. (Formally, $I(J)$ is the set of elements off
$J$ whose component of $G\setminus J$ has bounded point set.) An element $x\notin J$ is
\emph{within} $J$ if $x\in I(J)$ and \emph{without} $J$ if $x\in E(J)$. For an arc $pxq$ of elements,
$\oarc{pxq}=pxq\setminus\set{p,q}$.

\formal{\lean{HPC.Moore.sccInterior}, \lean{HPC.Moore.sccExterior}, \lean{HPC.Moore.sccInterior_spec},
\lean{HPC.Moore.closure_sccInterior}, \lean{HPC.Moore.closure_sccExterior}.}

\section{Arcs and curves of elements are topological arcs and curves}

The cross-cut theory of F.A.\ was developed in Chapter~\ref{ch:curves} for a space satisfying Moore's
axioms, but most of it uses the axioms only through a few consequences. Those consequences were isolated
in topological form: a Hausdorff, locally connected space $X$ with a family of open sets (``regions'')
such that
\begin{enumerate}[label=(\roman*)]
\item regions and the complements of their closures are connected, and region closures are compact,
while $X$ is not;
\item two points of a connected open set are joined by an arc in it (F.A.\ Theorem~15);
\item every simple closed curve is the frontier of exactly two components of its complement, one a region
(Axiom~8 together with F.A.\ Theorem~21).
\end{enumerate}
To apply this to $G$ we need Moore's arcs and curves of elements to be the topological ones.

\begin{proposition}\label{prop:elt-top}
\leavevmode
\begin{enumerate}[label=(\alph*)]
\item A set of elements is a simple continuous arc of elements from $A$ to $B$ if and only if it is the image of
an embedding of $[0,1]$ sending $0$, $1$ to $A$, $B$.
\item A set of elements is a simple closed curve of elements if and only if it is the union of two such arcs
with only their ends in common, if and only if it is homeomorphic to a circle.
\item Two arcs of elements from $p$ to $q$ with only $p,q$ in common form a simple closed curve of
elements.
\item $G$ is Hausdorff and arcwise connected; every connected open set of elements is arcwise connected.
\end{enumerate}
\end{proposition}

\begin{proof}
(a) An arc of elements is compact (Proposition~\ref{prop:arc-elt-top}), metrizable, connected, and every
element other than $A,B$ cuts it: by the cut-point characterization of the arc
(Theorem~\ref{thm:irr-arc} and Proposition~\ref{prop:cutarc}) it is a topological arc. Conversely a
topological arc is closed, connected, compact (so bounded: the point set of a compact set of elements is
compact, Corollary~\ref{cor:quot-tools}), and every interior point cuts it.

(b) A curve of elements is split by any two of its elements into two arcs of elements
(Theorem~\ref{thm:m25-5-10}(e)), which are topological arcs by (a). Conversely, two topological arcs with only
their ends in common are cut by any two points (Proposition~\ref{prop:scc-cut}). The circle characterization is
Theorem~\ref{thm:circle-char}.

(c) is (b). (d) Hausdorff: Theorem~\ref{thm:quot-metrizable}. Arcs: Moore's proof of F.A.\ Theorem~15 goes
through in $G$ with the Moore regions (Theorem~\ref{thm:fa-elementary} applied as in Chapter~\ref{ch:axioms});
alternatively, $G$ is a connected, locally connected, complete metric space, so connected open sets are
arcwise connected.
\end{proof}

\formal{\lean{HPC.Moore.isArcElt_iff_tarc}, \lean{HPC.Moore.isSCCElt_iff_tscc}, \lean{HPC.Moore.isSCCElt_union},
\lean{HPC.Moore.t2Space_quot}, \lean{HPC.Moore.arcConn_quot}, \lean{HPC.Jones.tscc_homeomorph_circle}.}

\begin{definition}\label{def:restricted}
A \emph{region in the restricted sense} is the interior $I(J)$ of some simple closed curve of elements $J$.
The \emph{space of elements with the restricted regions} is $G$ with these sets as Moore's regions.
\end{definition}

\formal{\lean{HPC.Moore.restrictedRegion}, \lean{HPC.Moore.spaceRestricted}.}

\begin{lemma}\label{lem:restricted-moore}
A region in the restricted sense is a Moore region (Definition~\ref{def:moore-region}). The restricted
regions satisfy (i)--(iii) above.
\end{lemma}

\begin{proof}
By Theorem~\ref{thm:m25-11}, $I(J)$ is open, connected and bounded with frontier $J$, which is connected. Its
closure $I(J)\cup J$ has bounded point set, so it is compact; the complement of the closure is $E(J)$, connected.
$G$ is not compact, since $\pts G=\R^2$ is unbounded. (ii) is Proposition~\ref{prop:elt-top}(d), and (iii) is
Theorem~\ref{thm:m25-11} together with Proposition~\ref{prop:elt-top}(b).
\end{proof}

\formal{\lean{HPC.Moore.mooreRegion_of_restricted}, \lean{HPC.Moore.th21Prop_restricted},
\lean{HPC.Moore.ax8Prop_restricted}, \lean{HPC.Moore.isCompact_closure_restricted},
\lean{HPC.Moore.not_isCompact_univ_quot}.}

Moore remarks (p.~423) that F.A.\ Theorem~21 ``may be easily proved with the help of Theorem~11''. This is
correct as it stands: only Theorem~\ref{thm:m25-11} is needed for (iii).

\section{Theorems 13 and 14}

\begin{lemma}[Janiszewski's theorem for elements]\label{lem:quot-jan}
Let $A,B$ be closed sets of elements with bounded point sets such that $A\cap B$ is connected (or empty). If
neither $A$ nor $B$ separates the elements $x$ and $y$, then $A\cup B$ does not separate them.
\end{lemma}

\begin{proof}
Components of $G\setminus A$ correspond to components of $\R^2\setminus\pts A$ through $H\mapsto\pts H$
(Proposition~\ref{prop:pts-conn}). The point sets $\pts A$, $\pts B$ are compact, $\pts A\cap\pts B=\pts{(A\cap B)}$
is connected (or empty), and Janiszewski's theorem in the plane (\vfour{Theorem~2.7}) applies to
points of $x$ and $y$.
\end{proof}

\formal{\lean{HPC.Moore.quot_janiszewski}.}

\begin{theorem}[Moore, Theorem~13]\label{thm:m25-13}
Let $R$ be a Moore region and $k$ either a single element or an arc of elements all of whose elements,
except possibly one end, belong to $R$. Then $R\setminus k$ is open and connected.
\end{theorem}

\begin{proof}
$R\setminus k$ is open since $k$ is closed. Let $x,y\in R\setminus k$ and $B=\fr R$. $B$ does not separate $x$
from $y$, since $R$ is connected and misses $B$. Nor does $k$: $G\setminus k$ is connected
(Theorem~\ref{thm:m25-9}, or Corollary~\ref{cor:pts-nonsep} for a single element). Now $k\cap B$ is empty or the one
end of $k$ outside $R$, so by Lemma~\ref{lem:quot-jan} $k\cup B$ does not separate $x$ from $y$. The component of
$G\setminus(k\cup B)$ containing $x$ and $y$ misses $B$ and meets $R$, so lies in $R$, and it misses $k$.
Finally $R\setminus k\ne\emptyset$ (a point not in Moore's text), since an arc of elements $k$ contains no nonempty open set of elements.
Indeed, if $U\subseteq k$ is open and nonempty, choose $y\in U$ other than the ends of $k$ and a Moore region $V\ni y$
inside $U$ (Theorem~\ref{thm:m25-4}). Then $V$ is a connected open subset of the arc $k$, so an interval of the arc containing $y$
in its interior, and $V\setminus\set y$ is not connected; but by the first part of this proof (with $V$ for $R$ and
$\set y$ for $k$) it is.
\end{proof}

\formal{\lean{HPC.Moore.not_subset_arc_of_cut}, \lean{HPC.Moore.m25_theorem13}.}

\begin{theorem}[Moore, Theorem~14]\label{thm:m25-14}
Every Moore region contains a simple closed curve of elements.
\end{theorem}

\begin{proof}
A Moore region $R$ is infinite, so it contains elements $p\ne q$, joined by an arc $pxq$ in $R$
(Proposition~\ref{prop:elt-top}(d)). Let $x$ be an element of $pxq$ other than $p,q$. By
Theorem~\ref{thm:m25-13}, $R\setminus\set x$ is a domain, so it contains an arc $pyq$ from $p$ to $q$. Moore:
``it is easy to see that the sum of the arcs $pxq$ and $pyq$ contains as a subset a simple closed curve''.
Indeed, follow $pxq$ from $x$ towards $p$ until the first element $a$ on $pyq$, and towards $q$ until the first
element $b$ on $pyq$ (both exist, since $p,q\in pyq$ and $pyq$ is closed). The sub-arc $axb$ of $pxq$ meets $pyq$ only
in $a$ and $b$, and with the sub-arc of $pyq$ from $a$ to $b$ it forms a simple closed curve of elements
(Proposition~\ref{prop:elt-top}(c)).
\end{proof}

\formal{\lean{HPC.Moore.exists_tscc_of_arcs}, \lean{HPC.Moore.m25_theorem14}.}

\section{Theorems 15--18: cross-cuts}

Moore gives no proofs of Theorems~15--18, only pointers to F.A.\ Theorems~24--27
(Theorems~\ref{thm:fa24}--\ref{thm:fa27}). Since those were proved from (i)--(iii) only, and
Lemma~\ref{lem:restricted-moore} supplies (i)--(iii) for $G$, each pointer is made exact: the F.A.\ proof, read in
$G$, with the axioms it uses replaced by Theorem~\ref{thm:m25-11}.

Fix arcs of elements $pxq$, $pyq$ with $pxq\cap pyq=\set{p,q}$, $J=pxq\cup pyq$, and a third arc $pzq$; put
$J_1=pxq\cup pzq$ and $J_2=pyq\cup pzq$.

\begin{theorem}[Moore, Theorems~15 and~16]\label{thm:m25-15-16}
Suppose $\oarc{pzq}\subseteq I(J)$. Then
\begin{enumerate}[label=(\arabic*)]
\item $I(J_1)\subseteq I(J)$;
\item $\oarc{pyq}\subseteq E(J_1)$;
\item $I(J_1)\cap I(J_2)=\emptyset$;
\item $I(J)=I(J_1)\cup I(J_2)\cup\oarc{pzq}$.
\end{enumerate}
\end{theorem}

\formal{\lean{HPC.Moore.m25_theorem15}, \lean{HPC.Moore.m25_theorem16}.}

\begin{theorem}[Moore, Theorem~17]\label{thm:m25-17}
If $pxq$, $pyq$, $pzq$ are three arcs of elements, no two with a common element other than $p$, $q$, then the
interiors of the three curves they form in pairs are not mutually disjoint.
\end{theorem}

\formal{\lean{HPC.Moore.m25_theorem17}.}

\begin{theorem}[Moore, Theorem~18]\label{thm:m25-18}
Suppose $\oarc{pzq}\subseteq E(J)$. Then (1) either $\oarc{pyq}\subseteq E(J_1)$ or $\oarc{pxq}\subseteq E(J_2)$; and
(2) if $\oarc{pyq}\subseteq E(J_1)$ then $\oarc{pxq}\subseteq I(J_2)$ and $I(J_2)=I(J)\cup I(J_1)\cup\oarc{pxq}$.
\end{theorem}

\formal{\lean{HPC.Moore.m25_theorem18}.}

Moore's statement of Theorem~18 uses the letters $x$, $y$ for elements of the arcs $pxq$, $pyq$ without
introducing them, and says ``without'' for ``in the exterior of''. We read the statements for all elements of the
open arcs. This reading loses nothing: $\oarc{pyq}$ is connected and misses $J_1$, so it lies wholly within or wholly
without $J_1$, and likewise for $\oarc{pxq}$ and $J_2$.

\section{Theorem 19 and the restricted topology}

\begin{theorem}[Moore, Theorem~19]\label{thm:m25-19}
If $R$ is a Moore region and $O\in R$, there is a simple closed curve of elements $J\subseteq R$ with
$O\in I(J)\subseteq R$.
\end{theorem}

\begin{proof}
Moore: ``with the use of Theorems 11, 12, 13, 17 and 18 this theorem may be proved by a method closely analogous to
that employed to prove Theorem 36 in F.A.'' We follow the proof of F.A.\ Theorem~36
(Theorem~\ref{thm:fa35-36}(b)), with its correction, in $G$, using Moore regions as the neighbourhood basis
(Theorem~\ref{thm:m25-4}).
\begin{enumerate}[label=(\arabic*)]
\item Choose an element $X_0\in R$, $X_0\ne O$, and a Moore region $Q\ni X_0$ in $R$ with $O\notin\cl Q$
(Theorem~\ref{thm:m25-3}). By Theorem~\ref{thm:m25-14}, $Q$ contains a closed curve $J$; its interior $I$
lies in $Q$ by Theorem~\ref{thm:m25-12} (applied to $D_1=Q$, $D_2=I$), so $O\in E(J)$.
\item An arc from $O$ to an element of $J$ inside $R$, cut at its first element on $J$, gives an arc $OA$
meeting $J$ only in $A$.
\item Since $R$ is bounded and $G$ is not, $\fr R\ne\emptyset$. By Theorem~\ref{thm:m25-13} (for $R\setminus OA$,
which is connected), there is an arc from an element of $J$ to a boundary element $C$ of $R$, missing $OA$ and,
except for $C$, in $R$; cut at its last element on $J\cup I$ (and its first element on $\fr R$) it gives an arc
$BC$ meeting $J\cup I$ only in $B$.
\item The elements $A,B$ split $J$ into arcs $AF_1B$ and $AH_1B$ (Theorem~\ref{thm:m25-5-10}(e)). By
Theorem~\ref{thm:m25-13} for $R$ minus the arc $OA\cup AF_1B\cup BC$ (an arc with one end, $C$, on $\fr R$),
there is an arc $FZH$ from $AF_1B$ to $AH_1B$, missing $OA\cup AF_1B\cup BC$ except for $F$, lying except for its ends in
$E(J)$ and in $R$.
\item $B$ is without the curve $ZHAFZ$: the interior of that curve lies in $R$ (Theorem~\ref{thm:m25-12}), while the arc
$BC$ reaches $C\notin R$ without meeting the curve. By Theorem~\ref{thm:m25-18}, $A$ is within $ZHBFZ$, and so is $O$,
since $OA$ does not meet that curve. Its interior lies in $R$ by Theorem~\ref{thm:m25-12}.
\end{enumerate}
\end{proof}

\formal{\lean{HPC.Moore.m25_theorem19}.}

Of Moore's citations, Theorems~11, 12, 13 and~18 are used; Theorem~17 is not (nor is F.A.\ Theorem~26 in the
proof of F.A.\ Theorem~36), while Theorems~10 and~14, which Moore does not cite, are.

\begin{theorem}[Moore, Theorem~20]\label{thm:m25-20}
An element $p$ is a limit element of a set $H$ of elements if and only if every region in the restricted sense
containing $p$ contains an element of $H$ other than $p$. More precisely: the restricted regions form a basis of the
topology of $G$, so the topology they generate is the quotient topology.
\end{theorem}

\begin{proof}
Restricted regions are open (Lemma~\ref{lem:restricted-moore}). Given an open $V\ni p$, take a Moore region
$R\ni p$ in $V$ (Theorem~\ref{thm:m25-4}) and then a curve $J$ with $p\in I(J)\subseteq R$ (Theorem~\ref{thm:m25-19}).
So the restricted regions form a basis, and limit elements (Theorem~\ref{thm:m25-4}) are the accumulation
points for this basis.
\end{proof}

\formal{\lean{HPC.Moore.isTopologicalBasis_restricted}, \lean{HPC.Moore.m25_theorem20},
\lean{HPC.Moore.spaceRestricted_topology}.}

\section{Theorems 21 and 22}

\begin{theorem}[Moore, Theorem~21]\label{thm:m25-21}
The space of elements with the restricted regions satisfies Moore's Axioms~1--8
(Definition~\ref{def:moore-axioms}).
\end{theorem}

\begin{proof}
By Theorem~\ref{thm:m25-20} and the bridge (Proposition~\ref{prop:bridge}), Moore's notions for this space (limit
points, closure, boundary, connectedness, arcs and closed curves) are the topological notions of $G$.

\emph{Axiom 1.} $G$ is a separable metric space (Theorem~\ref{thm:quot-metrizable}). Take a countable dense set
$\set{d_m}$ and, for each $m$ and $j$, a restricted region about $d_m$ inside the ball of radius $2^{-j}$ about
$d_m$ (Theorems~\ref{thm:m25-3} and~\ref{thm:m25-19}); enumerate these, each repeated infinitely often in a
suitable order. If $P\ne\bar P$ lie in a region $R$, the fundamental regions about $P$ of small radius have closures
in $R\setminus\set{\bar P}$.

\emph{Axioms 2, 4.} A restricted region is connected, and its closure compact (Lemma~\ref{lem:restricted-moore}).

\emph{Axiom 3.} The exterior of $I(J)$ is $G\setminus(I(J)\cup J)=E(J)$, connected (Theorem~\ref{thm:m25-11}).

\emph{Axiom 5.} Moore's limit elements are the same for both systems of regions (Theorems~\ref{thm:m25-4}
and~\ref{thm:m25-20}), so this is the corresponding statement for the Moore regions: the elements containing the
points $(n,0)$, $n$ large, form an infinite set without limit element.

\emph{Axiom 8.} A simple closed curve (topological, Proposition~\ref{prop:elt-top}(b)) is a curve of elements $J$, and
it is the boundary of the restricted region $I(J)$ (Theorem~\ref{thm:m25-11}).

\emph{Axioms 6, 7.} Moore derives these from F.A.\ and the fact that every restricted region is the interior of a
closed curve. We prove them directly. Let $R=I(J)$, $\bar R$ a restricted region and $P\in\bar R$ an element of $J$.
\begin{itemize}
\item \emph{Cross-cuts near $P$.} Let $V_0\subseteq\bar R$ be a Moore region about $P$ missing some element of $J$, and
$D=I(J)$ (or $D=E(J)$). There are elements $a,b$ of $J$ splitting $J$ into arcs $X$, $Y$ with $P\in\oarc X$ and
$X\subseteq V_0$, and an arc $c$ from $a$ to $b$ in $V_0$ with $\oarc c\subseteq D$. Indeed: the two arcs of $J$ from $P$
(Theorem~\ref{thm:m25-5-10}(e)), cut near $P$ (Theorem~\ref{thm:m25-5-10}), give $a$, $b$ on either side of $P$, which
can be taken accessible from $D$ inside a small region (an arc from an element of $D$ near $J$, cut at its first
element on $J$). The two access arcs are joined in $D\cap V_0$ by Lemma~\ref{lem:quot-jan}, applied to $J$ and to
$\fr V_0$ together with the rest of $J$; these meet in a connected set.
\item \emph{Axiom 6.} Take $c$ interior to $J$ and $K=I(X\cup c)$. By Theorem~\ref{thm:m25-15-16}, $K\subseteq R$, and by
Theorem~\ref{thm:m25-12}, $K\subseteq V_0\subseteq\bar R$. By Theorem~\ref{thm:m25-19} there is a restricted region $\bar K$
about $P$ inside $V_0\setminus Y$; then $J\cap\bar K\subseteq X$, which is part of the boundary of $K$.
\item \emph{Axiom 7.} Take $c$ exterior to $J$ and $L=I(X\cup c)$. By Theorem~\ref{thm:m25-18}, either $\oarc Y\subseteq
E(X\cup c)$, and then $L$ misses $J$ and lies in $E(J)$ (as $\oarc c$ does); or $\oarc Y\subseteq L\subseteq V_0$, which is
excluded since $V_0$ misses an element of $J$.
\end{itemize}
Along the way, Moore's weaker Axioms~$6'$ and~$7'$ (his route to 6 and 7 through F.A.) also follow, from
Theorems~\ref{thm:m25-11} and~\ref{thm:m25-13}.
\end{proof}

\formal{\lean{HPC.Moore.m25r_axiom1}, \lean{HPC.Moore.m25r_axiom2}, \lean{HPC.Moore.m25r_axiom3},
\lean{HPC.Moore.m25r_axiom4}, \lean{HPC.Moore.m25r_axiom5}, \lean{HPC.Moore.m25r_axiom6},
\lean{HPC.Moore.m25r_axiom7}, \lean{HPC.Moore.m25r_axiom8}, \lean{HPC.Moore.m25_axioms6'7'},
\lean{HPC.Moore.m25_theorem21}.}

\begin{theorem}[Moore, Theorem~22]\label{thm:m25-22}
If $G$ is a Moore decomposition of the plane, the space of elements $G$ is homeomorphic to $\R^2$. In Moore's form:
there is a one-to-one correspondence between the elements of $G$ and the points of the plane such that a point $P$ is
a limit point of a point set $M$ if and only if the corresponding element is a limit element of the set of
elements corresponding to $M$.
\end{theorem}

\begin{proof}
By Theorem~\ref{thm:m25-21} and Moore's characterization of the plane (Theorem~\ref{thm:moore-plane}), $G$ with the
topology generated by the restricted regions is homeomorphic to $\R^2$; by Theorem~\ref{thm:m25-20} this topology is the
quotient topology. Homeomorphisms preserve accumulation points, and limit elements are the accumulation points of
the quotient topology (Theorem~\ref{thm:m25-4}).
\end{proof}

\formal{\lean{HPC.Moore.m25_theorem22}, \lean{HPC.Moore.m25_theorem22_moore}.}

\section{Monotone decompositions of a plane continuum}\label{sec:monotone-quotient}

The main use of Theorem~\ref{thm:m25-22} in this series is to recognise a monotone image of a plane continuum as
a plane set.

\begin{proposition}\label{prop:monotone}
Let $M$ be a compact space and $q\colon M\to Y$ a continuous surjection onto a Hausdorff space with connected
fibres $q^{-1}(y)$. Then
\begin{enumerate}[label=(\alph*)]
\item the preimage of a connected set is connected;
\item if $M$ is an indecomposable continuum, so is $Y$ (provided $Y$ has more than one point): a decomposition
$Y=A\cup B$ into proper subcontinua lifts to the decomposition $M=q^{-1}(A)\cup q^{-1}(B)$.
\end{enumerate}
\end{proposition}

\begin{proof}
(a) $q$ is closed (compact to Hausdorff). If $q^{-1}(C)=U\cup V$ with $U,V$ disjoint, nonempty and relatively closed,
each fibre over $C$, being connected, lies in $U$ or in $V$, so $C=q(U)\cup q(V)$ is a partition of $C$ into disjoint,
nonempty, relatively closed sets. (b) follows from (a), as $q^{-1}(A)$ and $q^{-1}(B)$ are proper, closed and connected.
\end{proof}

\formal{\lean{HPC.Hagopian.isPreconnected_preimage_of_monotone},
\lean{HPC.Hagopian.isDecomposableSpace_of_monotone_image}.}

Now let $M\subseteq\R^2$ be a continuum and $q\colon M\to Y$ as above. The \emph{fibre decomposition} of the plane
has as elements the fibres of $q$ and the single points of $\R^2\setminus M$.

\begin{lemma}\label{lem:fibre-usc}
The fibre decomposition is upper semicontinuous in Moore's sense (Definition~\ref{def:moore-decomp}(3)).
\end{lemma}

\begin{proof}
For fibres of $q$: if fibres $F_n=q^{-1}(y_n)$ came within $1/n$ of $F=q^{-1}(y)$ but not within $e$, pass to a
subsequence with $y_n\to y'$ (compactness of $Y$) and with points $u_n\in F_n$ near $F$ converging to $u\in F$,
and points $w_n\in F_n$ far from $F$ converging to $w$. By continuity $q(u)=y'=q(w)$, so $y'=y$ and $w\in F$,
a contradiction. Points off $M$ are handled by the compactness of $M$: a point $z\notin M$ has a neighbourhood missing $M$,
whose elements are single points.
\end{proof}

\formal{\lean{HPC.Hagopian.fibreSetoid}, \lean{HPC.Hagopian.fibre_usc}.}

\begin{corollary}\label{cor:monotone-plane}
Let $M\subseteq\R^2$ be a continuum and $q\colon M\to Y$ a continuous surjection onto a compact Hausdorff space with
connected fibres, none of which separates the plane. Then the fibre decomposition is a Moore decomposition, the
space of elements is a plane, and $Y$ is homeomorphic to a compact subset of $\R^2$ (the image of $\pi(M)$).
In the application, $M$ is a homogeneous indecomposable plane continuum and every fibre is a proper subcontinuum of
$M$; such subcontinua do not separate the plane.
\end{corollary}

\begin{proof}
The elements are compact and connected and do not separate the plane by hypothesis (points never do), and the
decomposition is upper semicontinuous by Lemma~\ref{lem:fibre-usc}. By Theorem~\ref{thm:m25-22} there is a
homeomorphism $h\colon G\to\R^2$. The map $Y\to G$, $y\mapsto q^{-1}(y)$, is a continuous injection of a compact space
into a Hausdorff space (continuity: the fibre map is the quotient map restricted to $M$, factored through $q$, which is a
quotient map as a closed surjection), so it is an embedding; composing with $h$ embeds $Y$ onto the compact set
$h(\pi(M))$.
\end{proof}

\formal{\lean{HPC.Hagopian.fibreSetoid_mooreDecomp}, \lean{HPC.Hagopian.exists_plane_quotient_of_monotone_aux},
\lean{HPC.Hagopian.exists_plane_quotient_of_monotone}.}

The formal corollary is stated for the case of the application (homogeneous indecomposable $M$, fibres proper);
the non-separation of proper subcontinua is the input that lets Theorem~\ref{thm:m25-22} apply.

\section{Further plane tools}

Once the Schoenflies theorem is available in the plane itself, several classical facts used in later volumes
follow. We record them, with the plane-specific forms that are proved.

\begin{theorem}[Cross-cuts in the plane]\label{thm:plane-split}
Let $V\subseteq\R^2$ be a bounded connected open set whose frontier is the simple closed curve $C_1\cup C_2$, with
$C_1,C_2$ arcs from $a$ to $b$ meeting only in $a,b$. If $\gamma$ is an arc from $a$ to $b$ with $\oarc\gamma\subseteq V$, then
$V\setminus\gamma=R_1\cup R_2$, where $R_1,R_2$ are disjoint bounded connected open sets with frontiers
$\gamma\cup C_1$ and $\gamma\cup C_2$.
\end{theorem}

\begin{proof}
This is Theorems~\ref{thm:fa24} and~\ref{thm:fa25} for the restricted regions of the plane: the plane with the
interiors of its simple closed curves satisfies (i)--(iii) (Jordan curve theorem, Theorem~\ref{thm:jct}), and
$V=I(C_1\cup C_2)$.
\end{proof}

\formal{\lean{HPC.Hagopian.plane_split}.}

\begin{theorem}[Schoenflies theorem in the plane, one side prescribed]\label{thm:plane-schoenflies}
Let $J\subseteq\R^2$ be a simple closed curve, $D$ its bounded complementary domain, and $\varphi\colon[0,1]\to J$ an
embedding. There is an embedding $f\colon[0,1]^2\to\R^2$ with image $J\cup D$, $f(t,0)=\varphi(t)$ for all $t$, and $f$
maps the open square onto $D$.
\end{theorem}

\begin{proof}
Apply Theorem~\ref{thm:m25-21} to the decomposition of the plane into points (Section~\ref{sec:points}): the plane,
with the interiors of its simple closed curves as regions, satisfies Moore's Axioms~1--8. Cut the arc of $J$
complementary to $\varphi([0,1])$ twice to obtain a quadrilateral whose bottom side is $\varphi([0,1])$, apply
Theorem~\ref{thm:schoenflies} to it, and reparametrize the bottom side by $\varphi$.
\end{proof}

\formal{\lean{HPC.Hagopian.plane_schoenflies_aux}.}

\begin{theorem}[Two-sided collars]\label{thm:collar}
Every arc $\gamma$ in the plane has a two-sided collar: an embedding $F$ of $[0,1]\times[-1,1]$ onto a neighbourhood of
$\oarc\gamma$ (relative to $\gamma$'s complement near the ends) with $F(t,0)=\gamma(t)$.
\end{theorem}

\begin{proof}
Join the ends of $\gamma$ by an arc $\beta$ meeting $\gamma$ only in them (both ends are accessible from the connected
complement of $\gamma$). An arc $\beta_2$ from end to end through the bounded side $I$ of $\gamma\cup\beta$ cuts $I$ into the
bounded sides $R_1$ of $\gamma\cup\beta_2$ and $R_2$ of $\beta\cup\beta_2$ (Theorem~\ref{thm:plane-split}). One half-collar is
taken in $R_1\cup\gamma\cup\beta_2$ (Theorem~\ref{thm:plane-schoenflies}); the other in the unbounded side of $\gamma\cup\beta$,
by the Schoenflies theorem after an inversion about a point of $I$, restricted to a strip along $\gamma$.
\end{proof}

\formal{\lean{HPC.Hagopian.IsCollar}, \lean{HPC.Hagopian.exists_collar}.}

\section*{Exercises}
\addcontentsline{toc}{section}{Exercises}

\begin{exercise}
Show that in Theorem~\ref{thm:m25-13} the hypothesis that at most one end of the arc lies outside $R$ cannot be
dropped: give a Moore region $R$ and an arc with both ends on $\fr R$ and its other elements in $R$ such that $R\setminus k$
is not connected.
\end{exercise}

\begin{exercise}
Prove Theorem~\ref{thm:m25-17} for the decomposition of the plane into points, directly from the Jordan curve theorem.
\end{exercise}

\begin{exercise}
Let $M$ be the closed unit disc and $q\colon M\to[0,1]$, $q(z)=\abs z$. Describe the fibre decomposition. Which hypothesis of
Corollary~\ref{cor:monotone-plane} fails, and what is the space of elements?
\end{exercise}

\begin{exercise}
Let $G$ be the decomposition of $\R^2$ whose elements are the vertical segments $\set{x}\times[-1,1]$, $\abs x\le1$, and the
single points elsewhere. Verify the hypotheses of Theorem~\ref{thm:m25-22} and write down an explicit homeomorphism
of $G$ onto $\R^2$.
\end{exercise}
